% Author: Shuyu Fan, MSc
% Date: February 20, 2026

\subsubsection{Data Preprocessing Extensibility}
\label{subsubsec:preprocessing_extensibility}
The Harness knowledge base pipeline was originally designed around WHO SMART Guidelines distributed as FHIR (Fast Healthcare Interoperability Resources) artifacts. To support clinical content from other sources, we developed a parallel non-FHIR ingestion pathway that accepts common document formats and produces the same structured output used by the normalization and question generation pipelines.

\paragraph{Document Ingestion and Text Extraction}
The entry point for non-FHIR content is an extraction module that recursively scans an input directory and dispatches each file to a parser based on its extension.

\begin{verbatim}
Code: python/info_extraction/non_fhir_extraction.py
Function: extract_records_from_path
\end{verbatim}

Supported formats and their parsers:
\begin{itemize}
\item \textbf{PDF} -- Page-level text extraction via \texttt{pypdf}. Heading detection splits pages into topical sections before rechunking.
\item \textbf{DOCX} -- Paragraph and table extraction via \texttt{python-docx}. Table rows are flattened into pipe-delimited strings.
\item \textbf{HTML/HTM} -- Heading, paragraph, and list-item extraction via \texttt{BeautifulSoup}. Script, style, navigation, and footer elements are removed.
\item \textbf{TXT} -- Paragraph-level splitting on blank-line boundaries.
\item \textbf{JSON} -- A recursive string harvester that collects all string values from arbitrarily nested JSON structures, with deduplication.
\end{itemize}

After format-specific parsing, all extracted text passes through an overlapped rechunking step. The rechunker normalizes whitespace, produces chunks of up to 1,200 characters with 200-character overlap between consecutive chunks, and selects break points using a priority order of boundary characters (newline, period, semicolon, colon, comma, space). Because every format goes through the same rechunker, downstream components receive chunks of consistent size and structure.

Each chunk is represented as a \texttt{RawTextRecord} containing the source file path, chunk text, file extension, and per-file chunk index. The records are serialized to a JSONL file where each line contains:
\begin{verbatim}
{
  "source": "<file path>",
  "chunk_index": <global int>,
  "text": "<chunk text>",
  "meta": {"ext": ".pdf", "chunk_index_in_file": 0}
}
\end{verbatim}

\paragraph{LLM-Based Structured Extraction}
A ToolUniverse-backed extraction pipeline reads the chunked text and uses an LLM to identify diseases and extract structured clinical attributes.

\begin{verbatim}
Code: python/tool_universe/tooluniverse_parse_non_fhir.py
Function: main
\end{verbatim}

Chunks are grouped by source file and processed in configurable batches (default: 20 chunks per batch). For each batch, the combined text is passed to a \texttt{Harness\_Disease\_Extractor} AgenticTool configured with a structured prompt that instructs the LLM to return strict JSON conforming to the target schema:
\begin{verbatim}
{
  "<Disease Name>": {
    "importance": {"value": <number|null>,
                   "confidence": <0..1>},
    "symptoms":    {"values": [...], "confidence": <0..1>},
    "management":  {"values": [...], "confidence": <0..1>},
    "diagnosis":   {"values": [...], "confidence": <0..1>},
    "prognosis":   {"values": [...], "confidence": <0..1>},
    "contraindications": {"values": [...],
                          "confidence": <0..1>}
  }
}
\end{verbatim}

The extraction prompt requires concise noun phrases (3--4 words maximum) and case-insensitive deduplication. It also specifies a confidence scoring guide: 0.9--1.0 for information explicitly stated multiple times, 0.6--0.8 for information clearly stated once, 0.3--0.5 for implied evidence, and below 0.3 for weak or absent evidence.

The LLM provider is configurable via environment variables, supporting Gemini (default: Gemini 2.5 Pro) and OpenRouter. The LLM output is post-processed through schema normalization that clamps confidence scores to [0, 1], deduplicates and lowercases list values, and merges results across batches for the same disease.

\paragraph{Lineage Tracking}
The pipeline generates a separate lineage JSON file that maps each extracted value back to its source documents and chunk indices:
\begin{verbatim}
{
  "<Disease Name>": {
    "<field>": {
      "<extracted value>": [
        {"source": "<file path>",
         "chunk_indices": [0, 1, 3]}
      ]
    }
  }
}
\end{verbatim}

This lineage is carried through into the Neo4j knowledge graph, so every fact in the graph can be traced back to its original source document and specific text chunks.

\paragraph{Downstream Integration}
The non-FHIR extraction pipeline outputs a disease-centric JSON that is structurally identical to the FHIR-based pipeline's output. The normalization, Neo4j ingestion, and question generation components therefore work the same way whether the source content is a WHO SMART Guideline FHIR artifact or a PDF from an external clinical resource.

\paragraph{Adding New Document Formats}
To add a new format, a developer writes a parser function (e.g., \texttt{\_parse\_epub}) and registers it in the file dispatcher in \texttt{non\_fhir\_extraction.py}. The parser returns a list of text strings; the rechunker and all downstream components handle the rest.
